% titel: Blatt 6, A3: Polarkoordinaten
% skript: Kap. 1 (Polarkoordinaten, de Moivre)

% Teil I: Werkzeuge, kein Lösungsweg.

Zahlen in Polarkoordinaten umrechnen, Potenzen mit der Moivre-Formel
berechnen, komplexe Gleichungen wie $z^4=-1$ lösen und eine Drehung am
Ortsvektor durchführen.

\begin{werkzeugkasten}[Polarkoordinaten und de Moivre -- Potenzieren als Drehen]
Jede Zahl $z=a+ib$ lässt sich als Punkt lesen, der durch Abstand $r$ und
Winkel $\varphi$ bestimmt ist:
\begin{equation*}
  r=|z|=\sqrt{a^2+b^2},\quad
  \varphi=\arg z \ \text{im passenden Quadranten},\quad
  z=r(\cos\varphi+i\sin\varphi)=r\,\mathrm e^{i\varphi}.
\end{equation*}
Der Gewinn der Polardarstellung ist die \textbf{Multiplikation als
Drehung}: Mal $n$-mal dieselbe Zahl bedeutet $n$-mal dasselbe Drehen, also
Winkel $n\varphi$ und Radius $r^n$. Genau das besagt die
\textbf{Moivre-Formel}
\begin{equation*}
  z^n=r^n(\cos n\varphi+i\sin n\varphi)=r^n\,\mathrm e^{in\varphi}.
\end{equation*}
Ebenso entsteht die $n$-te Wurzel dadurch, dass man \emph{rückwärts dreht}
und alle $n$ Stellungen auf dem Kreis nimmt:
\begin{equation*}
  z^{1/n}=r^{1/n}\,\mathrm e^{i(\varphi+2\pi k)/n},\qquad k=0,1,\ldots,n-1 .
\end{equation*}
Die $+2\pi k$ stehen für die $n$ Winkel, die äquivalent zum selben Ort
sind; $z^n=w$ hat deshalb $n$ Lösungen.

\textit{Skript:} Kap. 1 (Polarkoordinaten, de Moivre).
\end{werkzeugkasten}

\begin{werkzeugkasten}[Kurze Umrechnungen -- Worauf zu achten ist]
\begin{itemize}
  \item $-5=-5+0i$: $r=5$, $\varphi=\pi$ (negative reelle Achse).
  \item $-3i+3=3-3i$: $r=\sqrt{9+9}=3\sqrt2$, $\varphi=-\tfrac\pi4$ (4.
        Quadrant).
  \item $\frac1i=\frac{i}{i^2}=-i$: $r=1$, $\varphi=-\tfrac\pi2$.
  \item $1+i\tan\alpha$: Realteil $\cos\alpha$, Imaginärteil $\sin\alpha$
        über die Skalierung $r=\sec\alpha=\frac1{\cos\alpha}$; es gilt
        $1+i\tan\alpha=\sec\alpha(\cos\alpha+i\sin\alpha)$, also
        $r=\sec\alpha$, $\varphi=\alpha$ (für $-\tfrac\pi2<\alpha<
        \tfrac\pi2$, wo $\cos\alpha>0$).
\end{itemize}
\textit{Skript:} Kap. 1 (Polarkoordinaten).
\end{werkzeugkasten}

\begin{achtung}[Typische Fallen]
\begin{itemize}
  \item \textbf{Quadrant des Arguments:} $1-i\sqrt3$ liegt im 4. Quadranten,
        $\varphi=-\tfrac\pi3$ (nicht $+\tfrac\pi3$). Den Winkel nie blind
        vom $\arctan$ übernehmen.
  \item \textbf{Bei $z^n$ IMMER erst in Polarkoordinaten umrechnen:}
        Potenzieren in $a+ib$-Form ist mühsam und fehleranfällig; de Moivre
        macht es zum Drehen.
  \item \textbf{$z^4=-1$ hat vier Lösungen} (nicht zwei): die vier Punkte
        $\mathrm e^{i\pi k/4}$ auf dem Einheitskreis.
  \item \textbf{Drehung um $\theta$ heißt Mal $\mathrm e^{i\theta}$:} der
        Radius bleibt, nur der Winkel $\varphi$ wird zu $\varphi+\theta$.
\end{itemize}
\end{achtung}
